\section{A Terse Review of Probability}

Both inequalities below are proved in a line, and between them they carry most of the arguments in these notes. The first needs nothing but the expectation.

\begin{boxtheorem}[Markov's Inequality] \label{Ch1:Thm:Markov}
    Let $X$ be a non-negative random variable and let $a > 0$. Then,
    \begin{align*}
        \Pr\of{X \geq a} \leq \frac{\E(X)}{a}
    \end{align*}
\end{boxtheorem}
\begin{proof}
    Compare $X$ with the random variable taking the value $a$ when $X \geq a$ and $0$ otherwise. As $X$ is non-negative, $X$ is everywhere at least as large, so its expectation is at least $a \Pr\of{X \geq a}$. Divide by $a$.
\end{proof}

The second inequality needs one more moment than the first.

\begin{boxdefinition}[Variance]
    The \textbf{variance} of a random variable $X$ is
    \begin{align*}
        \Var{X} = \E\of{\parenth{X - \E(X)}^2}
    \end{align*}
\end{boxdefinition}

Expanding the square and using linearity gives $\Var{X} = \E\of{X^2} - \E(X)^2$, which is usually the form to compute with. Expectation is linear whatever the variables do. Variance is not, but it does add over independent summands: for $X_1, \ldots, X_n$ independent,
\begin{align*}
    \Var{X_1 + \cdots + X_n} = \Var{X_1} + \cdots + \Var{X_n}
\end{align*}

% [CORRECTED] was: "Chebyshev's inequality says that for all $a \geq 0$" --- at $a = 0$ the left-hand side is
% $1$ and the right-hand side is $\Var{X} / 0$, so the statement there is not an inequality between real numbers
% at all; and both uses of it below take $a > 0$ anyway ($a = \E(X_n)$ and $a = \eps n$).
\begin{boxtheorem}[Chebyshev's Inequality] \label{Ch1:Thm:Chebyshev}
    Let $X$ be a random variable of finite variance and let $a > 0$. Then,
    \begin{align*}
        \Pr\of{\abs{X - \E(X)} \geq a} \leq \frac{\Var{X}}{a^2}
    \end{align*}
\end{boxtheorem}
\begin{proof}
    The event $\abs{X - \E(X)} \geq a$ is the event $\parenth{X - \E(X)}^2 \geq a^2$, and $\parenth{X - \E(X)}^2$ is non-negative with expectation $\Var{X}$. So \Cref{Ch1:Thm:Markov}, applied to it at $a^2$, is the bound.
\end{proof}

What gets used in practice is not Chebyshev itself but the following consequence of it, which turns smallness of the variance into the statement that $X$ is not zero.

\begin{boxcorollary} \label{Ch1:Cor:Second_Moment}
    Let $\parenth{X_n}_{n \in \N}$ be random variables with $\E(X_n) > 0$ for all large enough $n$. If $\frac{\Var{X_n}}{\E(X_n)^2} \to 0$, then $\Pr\of{X_n \neq 0} \to 1$.
\end{boxcorollary}
\begin{proof}
    If $X_n = 0$ then $\abs{X_n - \E(X_n)} = \E(X_n)$, so $\Pr\of{X_n = 0} \leq \Pr\of{\abs{X_n - \E(X_n)} \geq \E(X_n)}$, and \Cref{Ch1:Thm:Chebyshev} at $a = \E(X_n)$ bounds that by $\frac{\Var{X_n}}{\E(X_n)^2}$ for every such $n$, which tends to $0$.
\end{proof}

Chebyshev already says something worth seeing on its own.

% [CORRECTED] was: "$\Pr\of{\abs{X - \frac{n}{2}} \geq \eps_n}$" --- read the subscripted $\eps_n$ as the
% product $\eps n$, which is the only threshold at which Chebyshev returns the author's own right-hand side
% $\frac{\Var{X}}{\eps^2 n^2}$; at a genuine sequence $\eps_n$ the bound is $\frac{\Var{X}}{\eps_n^2}$ and the
% $\eps$ on the right is unbound.
\begin{boxexample}
    Flip $n$ fair coins and let $X$ be the number of heads. Then $\E(X) = \frac{n}{2}$, and as $X$ is a sum of $n$ independent variables of variance $\frac{1}{4}$ each, $\Var{X} = \frac{n}{4}$. So \Cref{Ch1:Thm:Chebyshev} at $a = \eps n$ gives
    \begin{align*}
        \Pr\of{\abs{X - \frac{n}{2}} \geq \eps n} \leq \frac{\Var{X}}{\eps^2 n^2} = \frac{1}{4 \eps^2 n}
    \end{align*}
    which tends to $0$ for every fixed $\eps > 0$: the proportion of heads is within $\eps$ of a half with probability tending to $1$.
\end{boxexample}
